\section{Decision (brief \S24) and next steps}
\paragraph{Outcome A (verified 3D reconstruction): not supported.} No independent reference (second view or 3D measurement) exists for any specimen. Input-view agreement after pose correction is a fit, not an accuracy measure.
\paragraph{Outcome B (knowledge-controlled 3D simulation): technical preconditions partly met; not established.} Three preconditions hold on the pilot:
\begin{itemize}
\item the KB-operator geometry editor gives coherent, regularised, KB-traceable shape change;
\item it realises more of the stage's KB geometric attributes than an equal-magnitude generic warp (small $n$, 2 classes, partly by construction);
\item one 3D asset is far more cross-view consistent than independent per-view 2D edits (internal consistency only).
\end{itemize}
Outcome B additionally requires \emph{biological validation}, which was not possible (no expert raters). The detectors are coarse; the depression detector is non-specific; on volumetric organs the outgrowth response is attenuated.
\paragraph{Outcome C (downstream augmentation): not tested.} The downstream classifier experiment (\S18) was not run. The Stage-5 gate requires a reliable controlled pilot first, and 15 assets cannot power a held-out comparison. At the measured $\approx$3.4 GPU-h per specimen for the full arm set, 250 assets need about 850 GPU-h locally. With only D, J and H per asset the figure is about 1 GPU-h per asset ($\approx$250 GPU-h, $\approx$5 days on two TITAN X).
\paragraph{Classification: Outcome D (no established advantage) for now.} Neither GaussianEditor nor the new pipeline has been shown to produce biologically valid disease geometry, and no advantage on real data has been measured. The evidence supports a narrower technical contribution: diagnosis and repair of the pipeline, plus a controllable, auditable KB$\to$geometry editing method with honest metrics.
\paragraph{Recommended next experiments, in priority order.}
\begin{enumerate}
\item \textbf{Blinded expert review} (2--3 plant pathologists). Rate stage-matched renders of D/J/H/K/F and 2D C-stage/I per specimen for symptom presence, stage correctness and implausible artefacts. This is the only way to move toward Outcome B.
\item \textbf{Operator coverage:} a directional \emph{roll} operator. Grape Leafroll's text specifies ``margins rolled under'', the clearest geometric KB claim in SAGE. Also calibrate outgrowth on synthetic volumes, and lesion-anchored (observed-spot) placement instead of random sites.
\item \textbf{Two-sided leaves:} SAM~3D leaves are single SH-0 sheets, so lower-surface symptoms (Tomato Leaf Mold, Squash Downy Mildew) cannot be represented. Duplicate the sheet with offset normals before any underside edit.
\item \textbf{Reference data for external validity:} even 20--30 specimens photographed from several views would turn every ``unverified'' label into a measurement. (Out of scope by current decision.)
\item \textbf{Downstream test (C)} only after 1--3: matched synthetic budget, leakage-safe SAGE test split, real-only vs +C-stage 2D vs +J/H 3D renders.
\end{enumerate}
